\chapter{Verbos irregulares básicos}
\section{Objetivos}
\begin{itemize}[leftmargin=*]
	\item Memorizar formas de presente de los verbos irregulares más frecuentes A1.
	\item Usarlos en preguntas y respuestas cortas.
\end{itemize}

\section{Contenidos clave}
\begin{itemize}[leftmargin=*]
	\item Formas completas de \textit{hebben}, \textit{zijn}, \textit{gaan}, \textit{komen}, \textit{doen}.
	\item Uso en afirmaciones, negaciones y preguntas con V2.
	\item Respuestas cortas con auxiliar + pronombre.
	\item Combinación con adverbios de lugar/tiempo básicos.
\end{itemize}

\section{Ruta del capítulo}
Seis secciones con conjugaciones y uso en frases.

\section{Verbos irregulares más frecuentes}
\textbf{Lista A1 clave:} \textit{zijn, hebben, gaan, komen, doen, willen, kunnen, moeten}.

\textbf{Por qué importan:} cubren la mayoría de interacciones cotidianas (estado, posesión, movimiento, obligación, capacidad, deseo).

\textbf{Práctica}
\begin{itemize}[leftmargin=*]
  \item Tarjetas: infinitivo + 3 formas esenciales; revisa a diario.
  \item Escucha diálogos e identifica cada aparición; repite la frase completa.
\end{itemize}

\section{\textit{hebben} en presente}
\textbf{Tabla breve}
\begin{center}
\begin{tabular}{lll}
\toprule
Persona & Forma & Ejemplo \\
\midrule
ik & heb & Ik heb een fiets. \\
je/jij & hebt & Heb je tijd? \\
u & heeft & Heeft u koffie? \\
hij/zij/het & heeft & Hij heeft een huis. \\
wij/jullie/zij & hebben & We hebben les. \\
\bottomrule
\end{tabular}
\end{center}

\textbf{Práctica}
\begin{itemize}[leftmargin=*]
  \item Crea 6 frases sobre posesiones/personas con \textit{hebben}.
  \item Pregunta-respuesta rápida en pareja: \textit{Heb je	extellipsis{}? Ja, ik heb	extellipsis{} / Nee, ik heb geen	extellipsis{}}.
\end{itemize}

\section{\textit{gaan} en presente}
\textbf{Tabla breve}
\begin{center}
\begin{tabular}{lll}
\toprule
Persona & Forma & Ejemplo \\
\midrule
ik & ga & Ik ga naar huis. \\
je/jij & gaat & Ga je mee? \\
u & gaat & Gaat u zitten. \\
hij/zij/het & gaat & Zij gaat morgen. \\
wij/jullie/zij & gaan & We gaan samen. \\
\bottomrule
\end{tabular}
\end{center}

\textbf{Práctica}
\begin{itemize}[leftmargin=*]
  \item 6 frases con destinos distintos (\textit{naar school, naar werk, naar de winkel}).
  \item Role-play: invitar a alguien con \textit{ga je mee?} y responder afirmando/negando.
\end{itemize}

\section{\textit{komen} en presente}
\textbf{Tabla breve}
\begin{center}
\begin{tabular}{lll}
\toprule
Persona & Forma & Ejemplo \\
\midrule
ik & kom & Ik kom uit Peru. \\
je/jij & komt & Kom je morgen? \\
u & komt & Komt u mee? \\
hij/zij/het & komt & Hij komt later. \\
wij/jullie/zij & komen & Ze komen samen. \\
\bottomrule
\end{tabular}
\end{center}

\textbf{Práctica}
\begin{itemize}[leftmargin=*]
  \item 6 frases de origen y llegada (\textit{Ik kom uit	extellipsis{}, We komen om acht uur}).
  \item Pregunta-respuesta con horarios: \textit{Hoe laat kom je?} / \textit{Ik kom om negen uur}.
\end{itemize}

\section{\textit{doen} en presente}
\textbf{Tabla breve}
\begin{center}
\begin{tabular}{lll}
\toprule
Persona & Forma & Ejemplo \\
\midrule
ik & doe & Ik doe sport. \\
je/jij & doet & Wat doe je morgen? \\
u & doet & Wat doet u? \\
hij/zij/het & doet & Zij doet yoga. \\
wij/jullie/zij & doen & We doen boodschappen. \\
\bottomrule
\end{tabular}
\end{center}

\textbf{Práctica}
\begin{itemize}[leftmargin=*]
  \item 6 frases con actividades diarias usando \textit{doen}.
  \item Role-play de pequeñas tareas: \textit{Doe je de afwas?} / \textit{Ja, ik doe het}.
\end{itemize}

\section{Uso en frases simples}
\textbf{Objetivo:} Integrar los irregulares en afirmaciones, negaciones y preguntas cortas.

\textbf{Plantillas}
\begin{itemize}[leftmargin=*]
  \item Afirmación: \textit{Ik heb tijd. We gaan naar huis. Zij komt straks.}
  \item Negación: \textit{Ik heb geen tijd. Hij komt niet. We doen het niet.}
  \item Pregunta: \textit{Heb je een minuut? Kom je mee? Wat doe jij?}
\end{itemize}

\textbf{Práctica}
\begin{itemize}[leftmargin=*]
  \item Escribe 10 frases mezclando 3 verbos irregulares diferentes.
  \item Haz un mini diálogo de 8 líneas usando \textit{hebben, gaan, komen, doen}.
\end{itemize}


\section{Vocabulario clave}
\begin{itemize}[leftmargin=*]
	\item Verbos ancla: \textit{hebben}, \textit{zijn}, \textit{gaan}, \textit{komen}, \textit{doen}.
	\item Adverbios frecuentes: \textit{nu}, \textit{straks}, \textit{hier}, \textit{daar}, \textit{morgen}.
\end{itemize}

\section{Pronunciación}
\begin{itemize}[leftmargin=*]
	\item Suaviza la /x/ en \textit{gaan} y evita pronunciar \textit{g} como /g/ español.
	\item \textit{hebben} pierde la \textit{n} final en habla rápida: \textipa{/'hE.b@/}.
\end{itemize}

\section{Errores comunes}
\begin{itemize}[leftmargin=*]
	\item Mezclar \textit{hebben} y \textit{zijn} en respuestas cortas; memoriza pares \textit{Ja, ik heb / Ja, ik ben}.
	\item Olvidar V2: \textit{Ga je morgen?} no \textit{Je gaat morgen?} en pregunta sí/no.
	\item Duplicar auxiliar en negaciones: usa solo \textit{doen} en énfasis, no con otros modales.
\end{itemize}

\section{Práctica guiada}
\begin{itemize}[leftmargin=*]
	\item Conjuga cada verbo en todas las personas con tarjetas.
	\item Transforma 10 frases afirmativas en preguntas y negaciones.
	\item Rellena respuestas cortas a preguntas dadas.
\end{itemize}

\section{Ejercicios de producción}
\begin{itemize}[leftmargin=*]
	\item Graba 10 frases hablando de planes (\textit{gaan}) y posesiones (\textit{hebben}).
	\item Escribe un mini diálogo de 8 turnos usando al menos 4 verbos irregulares distintos.
\end{itemize}

\section{Mini repaso}
\begin{itemize}[leftmargin=*]
	\item Conjugas los cinco irregulares sin mirar tablas.
	\item Colocas el verbo en segunda posición en afirmaciones y preguntas.
	\item Respondes de forma corta y natural (\textit{Ja, dat doe ik}).
\end{itemize}

\section{Resultado esperado}
\begin{itemize}[leftmargin=*]
	\item Al terminar puedes usar los irregulares básicos para hablar de posesión, movimiento y acciones diarias en afirmaciones, preguntas y negaciones cortas.
\end{itemize}

\section{Microevaluación}
\begin{itemize}[leftmargin=*]
	\item 10 preguntas/respuestas grabadas usando cada verbo al menos una vez.
\end{itemize}
